% Standalone empirical-results document for meanreversion_article_main.tex.
\iffalse
% Figure paths below are explicit and relative to this TeX file's project root.
% All numerical entries below are copied from files listed in Appendix B.
% Generated and theory-audited through the full-model gate on 2026-08-09.
\documentclass[11pt]{article}

\usepackage[T1]{fontenc}
\usepackage[utf8]{inputenc}
\usepackage{amsmath,amssymb,bm}
\usepackage{booktabs}
\usepackage{threeparttable}
\usepackage{graphicx}
\usepackage{subcaption}
\usepackage{geometry}
\usepackage{microtype}
\usepackage{url}
\usepackage[hidelinks]{hyperref}
\usepackage{xcolor}

\geometry{margin=1in}
\newcommand{\E}{\mathbb{E}}
\newcommand{\Var}{\operatorname{Var}}
\newcommand{\1}{\mathbf{1}}
\newcommand{\dd}{\,\mathrm{d}}
\newcommand{\kappahat}{\widehat{\kappa}}
\newcommand{\lambdahat}{\widehat{\lambda}}
\newcommand{\thetahat}{\widehat{\theta}}
\newcommand{\sourcefile}[1]{\par\noindent{\footnotesize\textit{Source file:} \path{#1}.}}

\title{Empirical Analysis of Mean Reversion in USDT/USD\\
\large Theory Alignment, Estimation, and Out-of-Sample Evidence}
\author{Empirical appendix for \texttt{meanreversion\_article\_main.tex}}
\date{9 August 2026}

\begin{document}
\maketitle

\begin{abstract}
This document implements and audits the stablecoin application of the copula-based
mean-reverting processes developed in \texttt{meanreversion\_article\_main.tex}.
The core parametric mixed-OU (MOU) and mixed-CIR (MCIR) transition kernels are
exactly the constant-intensity Poisson-switching special cases of the article's
mixed-copula construction.  The affine transformed OU and CIR specifications are
restricted versions of the article's linear transformation solutions.  The
nonparametric, threshold, rolling, and forecasting exercises are empirical
diagnostics that are not claimed as additional theoretical propositions.

Using 524,610 exact five-minute transition pairs from Kraken USDT/USD over
2021--2025, the estimated free-mean OU has a long-run log-price deviation of
0.665 basis points and a half-life of 589.5 minutes.  The jointly estimated MOU
substantially improves the in-sample transition likelihood.  A matched-design,
jointly re-estimated boundary bootstrap rejects the OU restriction
($LR=12{,}722.80$, $B=500$, corrected $p=0.001996$), although neither model passes
the residual specification tests and MOUF does not dominate OUF in the original
forecast exercise.  A diagnostic two-scale mixed copula attains the highest
out-of-sample copula log score, while the two-scale Gaussian copula most closely
matches the empirical dependence curve.  These two-scale specifications are not
one-dimensional Markov models and the empirical marginal is not verified to
satisfy the article's mean-reversion sufficient condition.  The evidence supports
the mixed kernel as a useful density extension, but does not identify its reset
intensity with a unique economic mechanism.  A second-stage registry estimates
or explicitly diagnoses all 25 finite-dimensional article and methodology
objects before ranking them.  On the untouched 2024 selection sample, an exact
filtered two-state Markov-switching OU has the highest five-minute conditional
log score (8.82469), ahead of a Student-$t$ marginal and Student-$t$ copula
(8.78120).  The day-block interval for the difference is strictly positive, but
both models have weakly identified structural parameters.  Model emphasis must
therefore distinguish predictive performance from structural identification.
\end{abstract}

\section{Relationship to the theoretical article}
\label{emp:sec:alignment}

The article defines a Markov process as mean reverting when it has an asymptotic
mean and its infinitesimal conditional drift has the sign of the displacement
toward that mean.  It then constructs new processes through (i) monotone
transformations of OU and CIR diffusions and (ii) mixtures of a consistent base
copula family with the independence copula.  Table~\ref{emp:tab:alignment}
records the exact relationship between those constructions and the empirical
implementation.

\begin{table}[htbp]
\centering
\caption{Theory-to-implementation alignment}
\label{emp:tab:alignment}
\begin{threeparttable}
\small
\begin{tabular}{p{0.22\textwidth}p{0.43\textwidth}p{0.25\textwidth}}
\toprule
Article object & Empirical implementation & Audit verdict \\
\midrule
Mean-reversion definition & Five-minute local-linear estimate of
$\mu_h(x)=h^{-1}\E[X_{t+h}-X_t\mid X_t=x]$ and sign statistic
$\1\{(\theta-x)\mu_h(x)>0\}$ & Finite-horizon proxy; it does not prove the
$h\downarrow0$ limit or asymptotic moment convergence. \\
Stationary OU and CIR & Exact OU transition density for $X_t=\log P_t$; CIR
transition density for $P_t$ and for absolute depeg pressure & Model equations
match.  Price CIR is numerically weakly identified. \\
Linear transformed OU &
$X_t=A_0e^{(\kappa-d)t}r_t+B_0e^{-dt}$ with $A_0=1$ and $B_0=0$ & Exact
restricted special case of the article's linear OU solution. \\
Linear transformed CIR &
$P_t=A_0e^{(\kappa-d)t}r_t+\theta[1-A_0e^{(\kappa-d)t}]+B_0e^{-dt}$,
with $A_0=1$ and $B_0=0$ & Exact restricted formula; base CIR parameters are
held fixed and the resulting model is weakly identified. \\
Mixed copula
$C^{\mathrm{mix}}_{s,t}=a_{s,t}C_{s,t}+(1-a_{s,t})\Pi$ & MOU and MCIR with
$a_{s,t}=e^{-\lambda(t-s)}$ and the stationary OU/Gamma marginal in the reset
component & Exact constant-intensity case of the article's Poisson construction. \\
General $a_{s,t}=\exp\{-\int_s^t c(u)\dd u\}$ & Constant $c(u)=\lambda$ only &
Time-varying reset intensity is not estimated. \\
Gaussian copula with a general marginal & Gaussian and mixed Gaussian copulas
with an empirical marginal & Copula construction is consistent, but the
article's sufficient sign condition for mean reversion is not imposed and is
not supported by the fitted finite-horizon drift. \\
Two-scale latent Gaussian copula & A weighted sum of fast and slow OU covariance
components, optionally combined with an independence reset & Diagnostic extension:
the two-dimensional latent state is Markov, but the observed scalar state is
generally not one-dimensional Markov and is not an article proposition. \\
Nonlinear OU/CIR transformations & Exponential, quadratic, and special-function
solutions are not estimated & Outside the present empirical scope. \\
Mixed-copula bond-pricing PDE & Not used for stablecoin prices & Not applicable
to the stated empirical objective. \\
\bottomrule
\end{tabular}
\begin{tablenotes}[flushleft]\footnotesize
\item The article's line following its diffusion equation states that
$\alpha(t)$ is a positive constant and then writes $\mu(x)=\alpha(x-\theta)$.
Under the article's own sign definition this is mean repulsion.  Consistency
requires either $\alpha<0$ or $\mu(x)=\alpha(\theta-x)$ for $\alpha>0$.  The
empirical code uses the latter, mean-reverting convention throughout.
\end{tablenotes}
\end{threeparttable}
\end{table}

The article does not specify an estimator in its currently empty empirical
section.  Maximum likelihood, composite likelihood, block bootstrap, PIT
diagnostics, and rolling out-of-sample scores are therefore empirical additions;
they do not alter the theoretical process definitions.

\section{Data, state variables, and exact transition design}
\label{emp:sec:data}

The raw file is Kraken's five-minute USDT/USD OHLCVT series.  Its SHA-256 digest
is
\begin{center}
\texttt{417ea3aa71f86c4476c30568f251f3f873ae8b0d44c5f480f5ad2f35fc6acc22}.
\end{center}
The audit finds 797,178 observations from 29 March 2017 to 31 December 2025,
750,220 exact 300-second adjacent pairs, no duplicate timestamps, no missing
fields, and grid coverage of 86.5186\%.  The main sample is 2021--2025.  The
training, validation, and test periods are 2021--2023, 2024, and 2025,
respectively.

\begin{table}[htbp]
\centering
\caption{Sample construction}
\label{emp:tab:samples}
\small
\begin{tabular}{lrrrr}
\toprule
Sample & Period & Observations & Coverage (\%) & Exact 5-minute pairs \\
\midrule
Baseline & 2021--2025 & 524,835 & 99.7998 & 524,610 \\
Extended & 2020--2025 & 627,782 & 99.4434 & 625,690 \\
Full & 2017--2025 & 797,178 & 86.5186 & 750,220 \\
Training & 2021--2023 & 314,759 & 99.8094 & 314,590 \\
Validation & 2024 & 105,153 & 99.7581 & 105,118 \\
Test & 2025 & 104,923 & 99.8126 & 104,900 \\
\bottomrule
\end{tabular}
\sourcefile{output/tables/sample_summary.csv}
\end{table}

For OU, threshold-OU, MOU, transformed-OU, and semiparametric copula models the
state is the log price
\begin{equation}
X_t=\log P_t,
\qquad X_t=0\Longleftrightarrow P_t=1.
\label{emp:eq:state}
\end{equation}
Price CIR and MCIR instead use $P_t>0$ directly.  Depeg-pressure CIR uses
$|P_t-1|+10^{-8}$ and is interpreted as an amplitude model, not a signed
mean-reversion model.  Transition pairs are formed by timestamp hashing: a pair
$(X_t,X_{t+h})$ is used only when the timestamp difference is exactly $h$.
Missing bars are never interpolated.

\section{Empirical specifications and estimators}
\label{emp:sec:methods}

\subsection{OU benchmark and finite-horizon drift}

The OU benchmark is
\begin{equation}
\dd X_t=\kappa(\theta-X_t)\dd t+\sigma\dd W_t,
\qquad \kappa>0,\quad\sigma>0.
\label{emp:eq:ou}
\end{equation}
For an interval $h$, its conditional mean and variance are
\begin{align}
m_h(x)&=\theta+e^{-\kappa h}(x-\theta),\\
v_h&=\frac{\sigma^2}{2\kappa}\left(1-e^{-2\kappa h}\right).
\label{emp:eq:oumoments}
\end{align}
OU0 fixes $\theta=0$, whereas OUF estimates $\theta$.  Parameters are estimated
by exact Gaussian transition maximum likelihood.  Ordinary Hessian standard
errors, UTC-day cluster-sandwich standard errors, and UTC-day block bootstrap
intervals are reported.

The model-free diagnostic estimates
\begin{equation}
\mu_h(x)=\frac{1}{h}\E[X_{t+h}-X_t\mid X_t=x]
\label{emp:eq:finite-drift}
\end{equation}
by local-linear regression on a 99-point quantile grid.  Bandwidth is selected by
blocked cross-validation and uncertainty is computed from 1,000 UTC-day bootstrap
replications.  The sign-consistency ratio is
\begin{equation}
\operatorname{SCR}_h(\theta)
=\frac{1}{G}\sum_{g=1}^G
\one\!\left\{(\theta-x_g)\widehat\mu_h(x_g)>0\right\}.
\label{emp:eq:scr}
\end{equation}

\subsection{Parametric mixed-copula transition kernel}

Let $P_h(x,\dd y)$ be the transition kernel of a stationary base process with
invariant distribution $F$.  The article's constant-intensity Poisson
construction gives
\begin{equation}
Q_h(x,\dd y)
=e^{-\lambda h}P_h(x,\dd y)
+\left(1-e^{-\lambda h}\right)F(\dd y).
\label{emp:eq:mixed-kernel}
\end{equation}
For an OU base process, the conditional density used in estimation is therefore
\begin{equation}
q_h(y\mid x)
=e^{-\lambda h}\phi\!\left(y;m_h(x),v_h\right)
+\left(1-e^{-\lambda h}\right)
\phi\!\left(y;\theta,\frac{\sigma^2}{2\kappa}\right).
\label{emp:eq:mou-density}
\end{equation}
This is exactly the transition density implied by the article's process
$Y_t=Y_t^{(N_t)}$ when the component processes are stationary OU processes.
It nests the OU at $\lambda=0$ and satisfies
\begin{equation}
\E[Y_{t+h}-\theta\mid Y_t=x]
=e^{-(\kappa+\lambda)h}(x-\theta).
\label{emp:eq:mou-mean}
\end{equation}
Hence the base diffusion half-life is $\log(2)/\kappa$, while the conditional
mean half-life of the mixed process is $\log(2)/(\kappa+\lambda)$.  These two
quantities must not be conflated.

MOU0 fixes $\theta=0$; MOUF jointly estimates
$(\theta,\kappa,\sigma,\lambda)$.  MCIR replaces the OU transition and Normal
invariant density in Equation~\eqref{emp:eq:mou-density} with the CIR transition
and its Gamma invariant density.  Because the price-CIR likelihood is weakly
identified, MCIR holds the base CIR parameters fixed and estimates only
$\lambda$; MCIR is therefore a restricted diagnostic, not a full joint MLE.

\subsection{Semiparametric multi-horizon copula}

Let $F_n$ be the empirical marginal CDF and $U_t=F_n(X_t)$.  For a Gaussian base
copula with $\rho_h=e^{-\kappa h}$, the fitted copula density is
\begin{equation}
c_h^{\mathrm{mix}}(u,v)
=e^{-\lambda h}c_{\mathrm{G}}(u,v;\rho_h)
+1-e^{-\lambda h}.
\label{emp:eq:semiparam-copula}
\end{equation}
The parameters are estimated by an equally weighted composite likelihood over
5, 10, 30, 60, 120, and 360 minutes.  Equation~\eqref{emp:eq:semiparam-copula}
defines a consistent stationary Markov copula construction.  It does
\emph{not}, for an arbitrary $F_n$, automatically satisfy the sign restriction
in the article's proposition for Gaussian copulas with general stationary
marginals.  This restriction is evaluated after estimation rather than assumed.

To test whether the failure at long horizons is caused by the single exponential
correlation curve, define the two-scale correlation
\begin{equation}
 \rho_h^{(2)}=w\exp(-\kappa_{f}h)+(1-w)\exp(-\kappa_{s}h),
 \qquad 0<w<1,\quad 0<\kappa_s<\kappa_f,
 \label{emp:eq:two-scale-rho}
\end{equation}
and the two-scale mixed density
\begin{equation}
 c_{h}^{(2,\mathrm{mix})}(u,v)
 =e^{-\lambda h}c_{\mathrm G}(u,v;\rho_h^{(2)})+1-e^{-\lambda h}.
 \label{emp:eq:two-scale-mix}
\end{equation}
Equation~\eqref{emp:eq:two-scale-rho} is generated by the sum of two independent
stationary OU factors.  The latent vector is Markov, but its scalar sum is in
general not one-dimensional Markov.  Consequently, Equations
\eqref{emp:eq:two-scale-rho}--\eqref{emp:eq:two-scale-mix} are used to diagnose
multi-horizon dependence and density forecasts; they are not presented as a new
instance of the article's one-dimensional Markov theorem.  Four models are fitted
on exactly the same deterministic subsample at each horizon: SingleGaussian,
SingleMixed, TwoScaleGaussian, and TwoScaleMixed.

\subsection{Prespecified model hierarchy}

The hierarchy is fixed before interpreting the additional results.  OUF and MOUF
are the \emph{core structural comparison}, because their transition laws directly
match the article.  Heteroskedastic threshold OU is the \emph{main empirical
extension} for conditional asymmetry.  The four empirical-marginal copula models
are \emph{diagnostic extensions}.  Random walk and nested-OU forecasts are
\emph{prediction benchmarks}.  CIR, MCIR, and affine transformed models are
\emph{appendix feasibility checks}.  A diagnostic model can reveal a missing
dependence scale without inheriting the theoretical status of the core models.
\sourcefile{output/tables/article_model_hierarchy.csv}

\subsection{Model comparison and inference}

Likelihood, AIC, and BIC are compared only for models fitted to the same state,
sample, and likelihood.  PIT uniformity and Ljung--Box tests assess transition
density calibration.  Because $\lambda=0$ is a boundary point, the ordinary
$\chi^2$ likelihood-ratio reference distribution is not used.  The primary
parametric bootstrap uses the same deterministic 25,000 pairs for the observed
and simulated statistics, holds the observed regressors fixed under the null, and
jointly re-estimates $(\theta,\kappa,\sigma,\lambda)$ under MOUF in every
replication.  Out-of-sample comparison uses negative
log score (NLS), CRPS, MAE, RMSE, Brier score, event log score, AUC, and calibration.
All 2025 predictions use parameters selected without 2025 outcomes.

\section{Empirical results}
\label{emp:sec:results}

\subsection{Nonparametric evidence and OU benchmark}

\begin{table}[htbp]
\centering
\caption{Finite-horizon mean-reversion sign diagnostics}
\label{emp:tab:scr}
\small
\begin{tabular}{lrrrr}
\toprule
Anchor definition & $\widehat\theta$ & SCR & 95\% bootstrap interval & Grid points \\
\midrule
Peg anchor & 0 & 0.7778 & [0.7677, 0.8788] & 99 \\
Sample anchor & $6.700\times10^{-5}$ & 0.8788 & [0.8586, 0.9697] & 99 \\
\bottomrule
\end{tabular}
\sourcefile{output/tables/sign_test_baseline.csv}
\end{table}

The point estimates predominantly have the correct direction, but the
simultaneous confidence band in Figure~\ref{emp:fig:drift} contains zero over a
large part of the state grid.  The evidence is therefore directional rather than
a proof of pointwise significant infinitesimal mean reversion.

\begin{figure}[htbp]
\centering
\includegraphics[width=0.82\textwidth]{output/figures/nonparametric_vs_ou_drift.pdf}
\caption{Local-linear finite-horizon drift and the OUF-implied drift.}
\label{emp:fig:drift}
\end{figure}

\begin{table}[htbp]
\centering
\caption{OU, asymmetric OU, and MOU parameter estimates}
\label{emp:tab:parameters}
\begin{threeparttable}
\small
\begin{tabular}{lrrrrrr}
\toprule
Model & $10^4\widehat\theta$ & $\widehat\kappa$ & $\widehat\lambda$ &
$\widehat\sigma$ & Base HL (min) & Effective HL (min) \\
\midrule
OUF & 0.665 & 1.6932 & --- & 0.001355 & 589.5 & 589.5 \\
MOUF & 7.248 & 0.02484 & 1.8863 & 0.000963 & 40,190.4 & 522.3 \\
\bottomrule
\end{tabular}
\begin{tablenotes}[flushleft]\footnotesize
\item Rates are per day.  OUF UTC-day cluster standard errors are
$(0.166,0.2099,0.0001456)$ for $(10^4\theta,\kappa,\sigma)$.  MOUF cluster
standard errors are $(6.496,0.01419,0.18835,0.00000950)$ for
$(10^4\theta,\kappa,\lambda,\sigma)$.  The large difference between the OUF and
MOUF estimates of $\theta$ is a warning against interpreting the MOUF stationary
mean as the economic peg without further stationarity restrictions.
\end{tablenotes}
\end{threeparttable}
\sourcefile{output/tables/model_estimates.csv}
\end{table}

The heteroskedastic threshold OU estimates
$(\kappa_+,\kappa_-)=(1.5674,1.8103)$ and
$(\sigma_+,\sigma_-)=(0.001248,0.001515)$.  Its likelihood gain is driven mainly
by asymmetric conditional volatility rather than by a large difference in
reversion speed.

\subsection{Does the mixed kernel add explanatory power?}

\begin{table}[htbp]
\centering
\caption{Comparable five-minute transition likelihoods}
\label{emp:tab:likelihood}
\small
\begin{tabular}{lrrrr}
\toprule
Model & Parameters & Log likelihood & AIC & BIC \\
\midrule
OU0 & 2 & 4,207,015.02 & $-8,414,026.04$ & $-8,414,003.70$ \\
OUF & 3 & 4,207,021.25 & $-8,414,036.50$ & $-8,414,002.98$ \\
Threshold OU, heteroskedastic & 4 & 4,211,719.53 & $-8,423,431.07$ & $-8,423,386.38$ \\
MOU0 & 3 & 4,351,515.34 & $-8,703,024.68$ & $-8,702,991.17$ \\
MOUF & 4 & 4,351,562.04 & $-8,703,116.07$ & $-8,703,071.39$ \\
\bottomrule
\end{tabular}
\sourcefile{output/tables/model_comparison.csv}
\end{table}

Relative to OUF, MOUF increases the total log likelihood by 144,540.79, or
approximately 0.2755 per transition pair, and lowers AIC by 289,079.58.  This is
the clearest in-sample evidence for the mixed kernel.  The gain reflects a
two-scale transition density: a highly persistent narrow base component and a
small-probability draw from a much wider invariant component.

The primary boundary bootstrap removes the design mismatch in the earlier
diagnostic.  Both the observed and null-simulated statistics use the same
deterministic systematic sample of 25,000 out of 524,610 exact transitions, the
same fixed regressor values and five-minute step, the same OUF refit, and the same
joint MOUF optimization.  The observed statistic is $LR=12{,}722.80$.  Among
$B=500$ null replications, none is as large (maximum $0.635$), giving the
finite-replication correction
\begin{equation}
 \widehat p=\frac{1+\sum_{b=1}^{B}\1\{LR_b\geq LR_{\mathrm{obs}}\}}{B+1}
 =\frac{1}{501}=0.001996.
 \label{emp:eq:bootstrap-p}
\end{equation}
The alternative optimizer reports convergence in 98.2\% of replications and
$\widehat\lambda$ is at its zero boundary in 99.4\%, as expected for this
nonregular null.  Thus the ordinary $\chi^2$ approximation is inappropriate, but
the matched conditional bootstrap provides strong evidence that the MOUF density
adds a component absent from OUF on this prespecified design.  The claim remains
conditional on the 25,000-pair systematic design and the maintained OU null; it is
not a proof that the economic data-generating process is MOUF.
\sourcefile{output/tables/matched_joint_mou_lr_bootstrap.csv}
\sourcefile{output/tables/matched_joint_mou_lr_bootstrap_draws.csv}

The former $B=200$ profiled calculation is retained only as a historical
sensitivity file because it compares a full-sample observed statistic with
shorter simulated paths and does not jointly fit the same alternative.  It is not
used for the primary significance statement.
\sourcefile{output/tables/mixed_lr_bootstrap.csv}

Three reported reset intensities answer different questions:
\begin{table}[htbp]
\centering
\caption{Reset-intensity estimates under different objectives}
\label{emp:tab:lambdas}
\small
\begin{tabular}{lrl}
\toprule
Estimation objective & $\widehat\lambda$ per day & Interpretation \\
\midrule
Joint five-minute MOUF likelihood & 1.8863 & Jointly changes the OU base and reset component \\
Matched joint MOUF likelihood & 1.9185 & Same 25,000-pair design as the boundary bootstrap \\
Restricted boundary profile & 0.7678 & Profiles $\lambda$ conditional on a re-estimated OU base \\
Multi-horizon empirical-marginal copula & 0.2557 & Composite dependence fit over 5--360 minutes \\
\bottomrule
\end{tabular}
\end{table}
These estimates are not interchangeable.  Their dispersion is direct evidence
that $\kappa$ and $\lambda$ are sensitive to the marginal model, estimation
objective, and horizon set.

\subsection{Multi-horizon dependence}

\begin{table}[htbp]
\centering
\caption{Empirical and fitted Gaussian-rank dependence}
\label{emp:tab:horizons}
\small
\begin{tabular}{rrrrrr}
\toprule
$h$ (min) & Pairs & Empirical & Base fit & Mixed fit & Reset weight \\
\midrule
5   & 524,610 & 0.99128 & 0.99692 & 0.99603 & 0.00089 \\
10  & 524,494 & 0.98868 & 0.99384 & 0.99208 & 0.00177 \\
30  & 524,279 & 0.98067 & 0.98165 & 0.97643 & 0.00531 \\
60  & 524,143 & 0.97160 & 0.96363 & 0.95342 & 0.01060 \\
120 & 524,017 & 0.95664 & 0.92858 & 0.90901 & 0.02108 \\
360 & 523,734 & 0.91630 & 0.80068 & 0.75111 & 0.06192 \\
\bottomrule
\end{tabular}
\sourcefile{output/tables/copula_horizon_fit.csv}
\end{table}

The mixed term reduces absolute dependence error at 5 and 10 minutes, but makes
the fit worse from 30 minutes onward.  At 360 minutes the data remain much more
persistent than either single-exponential specification.  A constant-intensity
independence reset therefore does not replace a slow latent factor, a regime
switching model, or a multi-scale OU structure.

The prespecified two-scale comparison confirms that the missing persistence is a
separate dimension from the reset mixture.
\begin{table}[htbp]
\centering
\caption{Single- versus two-scale empirical-marginal copula fits}
\label{emp:tab:scale-models}
\small
\begin{tabular}{lrrrrrr}
\toprule
Model & $\widehat\kappa_s$ & $\widehat\kappa_f$ & $\widehat w$ &
$\widehat\lambda$ & Mean log score & Dependence RMSE \\
\midrule
SingleGaussian   & 1.1398 & ---      & ---     & 0      & 1.42051 & 0.07022 \\
SingleMixed      & 0.8889 & ---      & ---     & 0.2505 & 1.46991 & 0.07018 \\
TwoScaleGaussian & 0.3348 & 251.0089 & 0.01244 & 0      & 1.52663 & 0.00377 \\
TwoScaleMixed    & 0.1361 & 225.0237 & 0.01215 & 0.5078 & 1.59503 & 0.03223 \\
\bottomrule
\end{tabular}
\begin{flushleft}\footnotesize
Rates are per day.  ``Mean log score'' is the equally weighted mean copula log
density over six horizons, using at most 100,000 deterministic pairs per horizon;
higher is better.  Dependence RMSE compares the six empirical Gaussian-rank
correlations with the fitted correlations; lower is better.  All four optimizers
converged, and repeated starts attained the same best solution.
\end{flushleft}
\sourcefile{output/tables/copula_scale_model_comparison.csv}
\end{table}

The two metrics answer different questions.  TwoScaleGaussian reduces dependence
RMSE by 94.6\% relative to SingleGaussian and reproduces the 360-minute dependence
as 0.9083 versus the empirical 0.9163.  TwoScaleMixed has the highest full-density
score, improving it by 0.0684 over TwoScaleGaussian, but its reset component lowers
the fitted 360-minute dependence to 0.8410 and therefore increases correlation
RMSE.  Thus the slow factor explains the long-memory shape, whereas the mixed
component explains additional features of the joint density; a likelihood gain
alone is not evidence of a better correlation curve.
\sourcefile{output/tables/copula_scale_horizon_fit.csv}

\begin{figure}[htbp]
\centering
\includegraphics[width=0.82\textwidth]{output/figures/single_vs_two_scale_dependence.pdf}
\caption{Empirical multi-horizon dependence and the four scale/reset specifications.}
\label{emp:fig:two-scale}
\end{figure}

\begin{figure}[htbp]
\centering
\begin{subfigure}{0.48\textwidth}
\centering
\includegraphics[width=\linewidth]{output/figures/gaussian_vs_mixed_copula.pdf}
\caption{Gaussian and mixed copula contours.}
\end{subfigure}\hfill
\begin{subfigure}{0.48\textwidth}
\centering
\includegraphics[width=\linewidth]{output/figures/multihorizon_dependence.pdf}
\caption{Dependence across horizons.}
\end{subfigure}
\caption{Copula shape and multi-horizon dependence.}
\label{emp:fig:copula}
\end{figure}

For the empirical marginal $F_n$, the fitted semiparametric conditional mean has
five-minute sign consistency of only 34.34\% around the exact peg and 42.42\%
around the empirical marginal mean $6.713\times10^{-5}$.  Thus the empirical-
marginal Gaussian/mixed-copula fit is useful as a dependence diagnostic, but it
is not empirically verified as a mean-reverting process under the article's
definition.  This conclusion is distinct from the parametric MOU result, whose
Normal invariant marginal and OU base satisfy the theoretical construction by
design.
\sourcefile{output/tables/copula_conditional_drift_baseline.csv}
\sourcefile{output/tables/descriptive_statistics_baseline.csv}

\subsection{Density diagnostics and out-of-sample performance}

Numerical optimization succeeds for OUF and MOUF, but both models are rejected
by PIT uniformity and by Ljung--Box tests of transformed residuals and squared
residuals (reported numerical $p$-values are effectively zero).  The mixed model
therefore improves relative likelihood without achieving absolute transition-
density adequacy.

\begin{figure}[htbp]
\centering
\includegraphics[width=0.86\textwidth]{output/figures/pit_and_residual_acf.pdf}
\caption{PIT and residual diagnostics for OUF and MOUF.}
\label{emp:fig:pit}
\end{figure}

\begin{table}[htbp]
\centering
\caption{Validation-period model selection, 2024}
\label{emp:tab:validation}
\small
\begin{tabular}{lrrrr}
\toprule
Model & Mean NLS & Mean CRPS & Mean RMSE & NLS rank \\
\midrule
OUF & $-7.18759$ & 0.00013091 & 0.00023667 & 1 \\
Random walk & $-7.13145$ & 0.00013279 & 0.00021416 & 2 \\
Nested-OU proxy & $-7.13006$ & 0.00014465 & 0.00026387 & 3 \\
MOUF & $-7.10424$ & 0.00015907 & 0.00029675 & 4 \\
\bottomrule
\end{tabular}
\sourcefile{output/tables/validation_model_selection.csv}
\end{table}

Lower NLS and CRPS are better.  OUF is selected on the validation period.  In
the 2025 test period, MOUF has a better 5-minute NLS for the 50-basis-point
depeg specification ($-8.7632$ versus $-8.7025$ for OUF), but OUF is better at
60, 360, and 1,440 minutes.  The mixed kernel therefore has a localized
short-horizon density advantage rather than uniform out-of-sample dominance.
\sourcefile{output/tables/forecast_scores.csv}

The scale/reset comparison is a separate copula-density exercise.  The empirical
marginal and all copula parameters are estimated using 2021--2023 only; the fitted
objects are then frozen for 2024 and 2025.  Table~\ref{emp:tab:copula-oos} reports
the equally weighted mean across the six horizons.
\begin{table}[htbp]
\centering
\caption{Out-of-sample copula log scores by model scale}
\label{emp:tab:copula-oos}
\small
\begin{tabular}{lrrr}
\toprule
Model & Training 2021--2023 & Validation 2024 & Test 2025 \\
\midrule
SingleGaussian   & 1.27385 & 1.54637 & 1.38939 \\
SingleMixed      & 1.31164 & 1.59861 & 1.43596 \\
TwoScaleGaussian & 1.40784 & 1.61073 & 1.50638 \\
TwoScaleMixed    & 1.47899 & 1.66640 & 1.57116 \\
\bottomrule
\end{tabular}
\begin{flushleft}\footnotesize
Higher log score is better.  Each period-specific mean first averages over
observations within a horizon and then gives equal weight to each of the six
horizons.  No validation or test outcomes enter parameter estimation.
\end{flushleft}
\sourcefile{output/tables/copula_scale_oos_scores.csv}
\end{table}

TwoScaleMixed exceeds SingleMixed by 0.06780 log-score units in validation and
0.13520 in test; it also exceeds TwoScaleGaussian by 0.05567 and 0.06478.  This is
the clearest out-of-sample metric in favor of combining a slow factor with the
mixed reset.  The advantage is not uniform by horizon: on the 2025 test data the
two-scale models score worse than SingleMixed at 5 and 10 minutes, but better at
30, 60, 120, and 360 minutes.  The result supports multi-scale dependence, not an
unqualified claim that the mixed process dominates every forecast horizon.

The baseline 50/25-basis-point event rule identifies eight events.  Their
empirical recovery-time median is 155 minutes, compared with 665 minutes under
OUF simulation and 1,085 minutes under MOUF simulation.  The reset mechanism
does not necessarily move the state toward the peg; it draws from the invariant
marginal.  Consequently, the effective conditional-mean half-life in
Equation~\eqref{emp:eq:mou-mean} is not a first-passage-time statistic.
\sourcefile{output/tables/first_passage_comparison.csv}

\subsection{Frequency, sample, and block-length robustness}

\begin{table}[htbp]
\centering
\caption{Baseline frequency robustness}
\label{emp:tab:frequency}
\small
\begin{tabular}{rrrrrr}
\toprule
Minutes & Exact pairs & $\widehat\kappa$ & $\widehat\lambda$ & OU HL & Effective HL \\
\midrule
5  & 524,610 & 1.6932 & 0.6969 & 589.5 & 417.6 \\
15 & 174,578 & 1.8504 & 0.1784 & 539.4 & 492.0 \\
30 & 87,068  & 1.3992 & 0.1323 & 713.4 & 651.7 \\
60 & 43,322  & 1.1311 & 0.0785 & 882.5 & 825.2 \\
\bottomrule
\end{tabular}
\begin{flushleft}\footnotesize
All non-five-minute rows use complete aggregation bins only.  Half-lives are in
minutes.  In this table $\lambda$ is estimated conditionally on fast OU
parameters, so it should not be compared numerically with the joint MOUF estimate
without qualification.
\end{flushleft}
\sourcefile{output/tables/robustness_estimates.csv}
\end{table}

The decline in $\widehat\kappa$ at lower sampling frequencies is consistent with
microstructure noise inflating high-frequency reversion estimates.  The 1-, 3-,
and 7-day moving-block bootstrap standard errors for OUF $\kappa$ are 0.2021,
0.2300, and 0.2622, respectively; uncertainty increases with block length but
the reversion-rate estimate remains positive.  In the full 2017--2025 sample,
retaining three flagged hard anomalies produces $\widehat\kappa=303.09$, whereas
excluding them gives $\widehat\kappa=2.742$.  Structural interpretation of the
full-sample high-frequency estimate is therefore untenable without explicit
outlier treatment.
\sourcefile{output/tables/ou_block_length_robustness.csv}

\section{Full-model registry and validation selection}
\label{emp:sec:full-registry}

The second-stage exercise removes the MVP complexity screen.  Every article
model with an explicit finite-dimensional observation density, and every
benchmark or marginal--copula family required by the empirical methodology, is
either estimated or assigned a terminal failure diagnostic before any model is
chosen for emphasis.  The registry contains 25 theory/methodology objects.  Of
these, 18 are estimated in the current or pre-existing pipeline, five terminate
with a retained monotonicity, support, or endpoint failure, the stationary
Gaussian--Gamma construction is estimable but fails the article's shape
condition, and the bond-pricing PDE is retained as not observable from spot data.
\sourcefile{output/tables/full_model_implementation_status.csv}

For every comparable model, selection uses the same 2024 five-minute
conditional log score on the log-price measure,
\begin{equation}
 \widehat S_m^{V}
 =\frac{1}{N_V}\sum_{i=1}^{N_V}\log q_m
   (x_{t_i}\mid\mathcal F_{t_i-}),
 \qquad N_V=40{,}000.
 \label{emp:eq:uniform-score}
\end{equation}
Price-density models are converted by the Jacobian
$q_X(x_t\mid x_{t-h})=q_P(e^{x_t}\mid e^{x_{t-h}})e^{x_t}$.
Copula models use
$q_X(x_t\mid x_{t-h})=c_h(F(x_{t-h}),F(x_t))f(x_t)$.
Thus AIC values from ordinary and composite likelihoods do not enter this
cross-class ranking.  For the two-state OU, the primary one-step score is the
exact hidden-state predictive density
\begin{equation}
 q_t(x_t\mid\mathcal F_{t-1})
 =\sum_{j=1}^{2}\left(\sum_{i=1}^{2}\alpha_{t-1,i}p_{ij}\right)
 p_j(x_t\mid x_{t-1}),
 \label{emp:eq:msou-filter-score}
\end{equation}
with the filter reset to its stationary probabilities after every timestamp
gap.  Longer-horizon Markov-switching entries are labelled as endpoint-mixture
approximations and do not determine selection.

\begin{table}[htbp]
\centering
\caption{Unified 2024 validation ranking at the five-minute horizon}
\label{emp:tab:full-ranking}
\begin{threeparttable}
\small
\begin{tabular}{rlrrl}
\toprule
Rank & Model & Mean log score & Coverage & Identification audit \\
\midrule
1 & Markov-switching OU & 8.824686 & 1.000 & weak: slow regime/boundary \\
2 & Student-$t$ margin + Student-$t$ copula & 8.781202 & 1.000 & weak: $\nu=2.05$ grid boundary \\
3 & Student-$t$ margin + mixed Student-$t$ copula & 8.781201 & 1.000 & weak: $\lambda\simeq0$, $\nu$ boundary \\
23 & MOUF & 8.431420 & 1.000 & passes numerical identification gate \\
43 & OUF & 8.221318 & 1.000 & passes numerical identification gate \\
44 & Random walk & 8.220764 & 1.000 & no-reversion benchmark \\
\bottomrule
\end{tabular}
\begin{tablenotes}[flushleft]\footnotesize
\item The rank numbers for MOUF, OUF, and the random walk are their positions in
the complete candidate table; omitted rows are not omitted from estimation or
ranking.  Selection uses 2024 only.  The 2025 scores are post-extension
exploratory because aggregate 2025 outcomes had been inspected previously.
\end{tablenotes}
\end{threeparttable}
\sourcefile{output/tables/full_model_validation_ranking.csv}
\end{table}

The selected model's advantage is not merely a point-ranking artifact.
Day-block resampling of paired score differences gives the comparisons in
Table~\ref{emp:tab:full-pairwise}.  These intervals measure predictive-score
uncertainty conditional on the frozen estimates; they do not repair weak
identification of the switching parameters.

\begin{table}[htbp]
\centering
\caption{Paired day-block comparison of the selected model}
\label{emp:tab:full-pairwise}
\small
\begin{tabular}{lrrr}
\toprule
Reference & $\widehat S_{MSOU}-\widehat S_{ref}$ & 95\% block interval & Daily win share \\
\midrule
Student-$t$ margin + Student-$t$ copula & 0.043484 & [0.018242, 0.065954] & 0.656 \\
MOUF & 0.393266 & [0.355106, 0.429977] & 0.915 \\
OUF & 0.603368 & [0.566379, 0.640416] & 0.956 \\
Random walk & 0.603922 & [0.564066, 0.642485] & 0.959 \\
\bottomrule
\end{tabular}
\sourcefile{output/tables/full_model_pairwise_validation.csv}
\end{table}

The expanded grid also makes the mixed-copula claim more precise.  With the
Student-$t$ marginal fixed, adding the independent-reset component raises the
Gaussian validation score by 0.078347 and the two-scale Gaussian score by
0.048432.  It changes the Student-$t$ copula score by only
$-6.95\times10^{-7}$ because $\widehat\lambda$ is effectively zero.  For Frank,
Clayton, survival Clayton, Gumbel, and survival Gumbel, the fitted mixed versions
hit weak-identification/boundary diagnostics and reduce the validation score by
1.98 to 3.21.  The empirical advantage is therefore specific to the Gaussian
and two-scale Gaussian reset constructions; it is not a theorem that mixing
improves every copula family.
\sourcefile{output/tables/full_mixture_incremental_scores.csv}

The transformation audit is equally informative.  Affine OU is estimable;
exponential OU converges but is weakly identified.  The OU special-function
solution is even in the latent state and hence non-injective on $\mathbb R$.
Affine CIR is estimated but weak, while quadratic CIR, exponential CIR, and the
prespecified CIR special-function branches fail global monotonicity or observed
support.  These failures are substantive outcomes, not silent exclusions.  The
Gamma Case~1 equation has no finite shape root on the audited range, so Case~2
lacks the stationary endpoint it requires.  Finally, the mixed-copula bond PDE
cannot be estimated without a bond-price cross section and a risk-neutral
measure; its equation is implemented only as a non-observable theoretical
object in this spot-price application.

\section{What the evidence establishes}
\label{emp:sec:conclusion}

The empirical evidence supports the following ordered conclusions.
\begin{enumerate}
\item USDT/USD log prices display finite-horizon directional reversion around a
near-zero log-price anchor.  A simple OU provides a transparent first benchmark,
but its transition density is misspecified.
\item The article's parametric mixed-copula construction is implemented exactly
for stationary OU and CIR bases.  For the OU base it yields a very large
in-sample likelihood gain.  The same-design, jointly re-estimated boundary
bootstrap rejects the OU restriction at the minimum attainable corrected level
$1/501$, conditional on the prespecified 25,000-pair design.
\item The gain should be interpreted as evidence for a two-component transition
density, not as causal evidence that $\lambda$ measures redemption, arbitrage,
news, or any unique economic event.
\item MOUF does not dominate OUF in the original validation exercise, PIT
diagnostics, or empirical first-passage times.  A constant independent-reset
intensity is therefore an incomplete empirical description of the price process.
\item Adding a slow latent factor reduces six-horizon dependence RMSE from 0.07022
to 0.00377.  Adding the reset mixture to the two-scale model raises validation and
test copula log scores to 1.66640 and 1.57116, the best of the four diagnostic
models, but worsens correlation RMSE to 0.03223.  Persistence shape and density
fit are distinct empirical targets.
\item The empirical-marginal semiparametric copula is not shown to satisfy the
article's sufficient mean-reversion condition.  It must be labelled as a
dependence and forecasting extension unless the marginal sign restriction is
imposed during estimation.
\item All article transformation families with explicit finite-dimensional
forms have now been implemented.  Their failures matter empirically: nonlinear
CIR variants fail monotonicity/support, exponential OU is weakly identified,
and the OU special-function map is non-injective.  The affine-OU calibration is
the only transformation result suitable for routine comparison.
\item Complexity-blind 2024 validation selects the two-state Markov-switching OU
on predictive log score, not a mixed-copula model.  Its paired advantage over the
best Student-$t$ copula model is positive, but its near-zero slow-regime reversion
rate is weakly identified.  It may be emphasized as a predictive benchmark,
while structural interpretation of its regimes must remain qualified.
\end{enumerate}

The most defensible statement of the mixed-copula advantage is therefore:
\emph{relative to an OU transition, the Poisson-switching mixed kernel captures a
narrow persistent component and rare wide transitions, produces a large matched-
design in-sample likelihood gain, and preserves the base invariant marginal and
Markov semigroup.  Within the diagnostic two-scale family, the mixed reset also
improves validation and test copula log scores for Gaussian and two-scale
Gaussian dependence.}  Its limitation is equally clear: the reset does not by
itself generate the observed slow persistence, is redundant for the fitted
Student-$t$ copula, fails out of sample for the fitted Archimedean mixed variants,
worsens the two-scale correlation RMSE, and does not reproduce empirical
recovery times.  It remains the article's main theoretical contribution, but the
empirical headline model is chosen by the frozen validation rule rather than by
ownership or complexity.

\fi
\clearpage
\section{Theory and implementation audit checklist}
\label{emp:app:audit}

\begin{enumerate}
\item[\checkmark] MOU density equals
$e^{-\lambda h}p_{\mathrm{OU}}+(1-e^{-\lambda h})f_{\mathrm{OU,stat}}$ pointwise.
The code test verifies exact equality to OU when $\lambda=0$.
\item[\checkmark] MCIR uses the stationary Gamma law implied by the same
$(\theta,\kappa,\sigma)$ as the CIR base.
\item[\checkmark] Mixed conditional mean decays at $\kappa+\lambda$; simulation
and analytic moments are compared in the test suite and generator-validation
table.
\item[\checkmark] Affine OU and CIR inverse transformations and Jacobians match
the article's linear formulas under $A_0=1$ and $B_0=0$.
\item[\checkmark] All primary five-minute pairs are exact timestamp matches; no
missing bar is treated as an observed transition.
\item[$\triangle$] The article assumes a stationary base process.  Rolling
estimates show substantial parameter variation, so stationarity is an empirical
approximation rather than an established fact.
\item[$\triangle$] The semiparametric empirical marginal does not impose the
article's Gaussian-copula marginal sign restriction and fails the finite-horizon
sign diagnostic on most grid points.
\item[$\triangle$] Price-CIR likelihood evaluation requires a high-noncentrality
Gaussian moment fallback for all evaluated price transitions; price CIR and
affine transformed CIR are flagged as weakly identified.
\item[\checkmark] The primary boundary bootstrap uses the same 25,000 transition
pairs and the same joint OUF/MOUF estimators for the observed and all 500 simulated
statistics.  Its inference is conditional on the systematic design and fixed
observed regressors, rather than an unconditional full-sample statement.
\item[$\triangle$] The two-scale scalar copula model captures multi-horizon
dependence but is not one-dimensional Markov; it is a diagnostic latent-factor
extension rather than an implementation of the article's scalar Markov theorem.
\item[\checkmark] Deterministic seasonal $c(t)$, exponential OU, quadratic and
exponential CIR, and the article's special-function branches are implemented.
Converged, weak, non-monotone, and support-failed outcomes are all retained in
the same implementation-status table.
\item[\checkmark] Normal, Student-$t$, skew Student-$t$, GED, and monotone-spline
margins are crossed with Gaussian, Student-$t$, Frank, Clayton, survival Clayton,
Gumbel, and survival Gumbel copulas; base and independence-mixed versions are
scored for all three periods and four horizons.
\item[\checkmark] Random-walk, jump-OU, nested-OU, two-state OU, and a
finite-dimensional deterministic time-varying marginal--copula model enter the
same primary validation ranking.  No model is removed because it is complex.
\item[$\triangle$] Markov-switching OU wins the exact filtered one-step score but
its slow-regime rate is near the optimization boundary and spans substantially
less than one mean-reversion unit in the continuous estimation segment.
\item[$\times$] The mixed-copula bond-pricing PDE is not empirically observable
from a USDT spot-price series without a bond-price cross section and a specified
risk-neutral measure.  This is retained as an explicit data-identification
failure rather than an unimplemented model.
\end{enumerate}

\section{Output provenance}
\label{emp:app:provenance}

\small
\begin{longtable}{p{0.42\textwidth}p{0.50\textwidth}}
\caption{Primary reproducibility files}\label{emp:tab:provenance}\\
\toprule
File & Content \\
\midrule
\endfirsthead
\toprule
File & Content \\
\midrule
\endhead
\path{output/tables/data_audit.csv} & Raw-data hash, rows, timestamps, gaps, and audit conditions \\
\path{output/tables/sample_summary.csv} & Sample periods, coverage, and exact-pair counts \\
\path{output/tables/sign_test_baseline.csv} & Nonparametric sign-consistency statistics \\
\path{output/tables/model_estimates.csv} & Parameters and ordinary, cluster, and bootstrap uncertainty \\
\path{output/tables/model_comparison.csv} & Likelihood, information criteria, numerical checks, and PIT tests \\
\path{output/tables/matched_joint_mou_lr_bootstrap.csv} & Primary matched-design joint-MOUF boundary bootstrap \\
\path{output/tables/matched_joint_mou_lr_bootstrap_draws.csv} & All 500 null LR draws and convergence diagnostics \\
\path{output/tables/mixed_lr_bootstrap.csv} & Superseded restricted boundary-bootstrap sensitivity \\
\path{output/tables/copula_horizon_fit.csv} & Empirical and fitted dependence over six horizons \\
\path{output/tables/copula_scale_model_comparison.csv} & Single/two-scale fit, parameters, and dependence RMSE \\
\path{output/tables/copula_scale_horizon_fit.csv} & Four-model fitted dependence at each horizon \\
\path{output/tables/copula_scale_oos_scores.csv} & Frozen-training-model validation and test copula log scores \\
\path{output/tables/article_model_hierarchy.csv} & Structural, extension, diagnostic, benchmark, and appendix tiers \\
\path{output/tables/copula_conditional_drift_baseline.csv} & Semiparametric conditional means and sign diagnostic \\
\path{output/tables/forecast_scores.csv} & Validation and test forecast scores \\
\path{output/tables/first_passage_comparison.csv} & Empirical and simulated recovery-time distributions \\
\path{output/tables/robustness_estimates.csv} & Frequency, extended-sample, and anomaly-policy robustness \\
\path{output/tables/full_model_registry.csv} & Frozen article/methodology model universe and comparison groups \\
\path{output/tables/full_model_implementation_status.csv} & Convergence, weak identification, monotonicity, support, and retained failures \\
\path{output/tables/full_model_validation_ranking.csv} & Complexity-blind 2024 five-minute conditional-density ranking \\
\path{output/tables/full_model_pairwise_validation.csv} & Paired day-block uncertainty for the selected model \\
\path{output/tables/full_mixture_incremental_scores.csv} & Within-family mixed-minus-base score increments \\
\path{output/tables/updated_stage_gate.csv} & Full-model coverage, scoring, diagnostics, and no-leakage gates \\
\path{output/models/mixed_fits_baseline.json} & Full MOU/MCIR parameter and optimization records \\
\path{output/models/full_model_fits.json} & Full-suite fits and nonlinear failure diagnostics \\
\path{output/models/full_copula_fits.json} & Five margins and fourteen fitted copula specifications \\
\path{output/models/semiparametric_copula_baseline.json} & Multi-horizon composite-likelihood estimate \\
\path{output/models/copula_scale_models.json} & Full-sample and training-sample scale/reset estimates \\
\path{output/run_status.json} & Step-level one-click pipeline status \\
\path{output/manifest.json} & Artifact hashes and joint test/pipeline completion status \\
\bottomrule
\end{longtable}
\normalsize

\endinput
